% Fork della blockchain di Bitcoin: evoluzione della rete in cinque passi.
% Il colore di ogni nodo indica l'ultimo blocco della catena che il nodo considera valida.
% Incluso nel documento con % Fork della blockchain di Bitcoin: evoluzione della rete in cinque passi.
% Il colore di ogni nodo indica l'ultimo blocco della catena che il nodo considera valida.
% Incluso nel documento con % Fork della blockchain di Bitcoin: evoluzione della rete in cinque passi.
% Il colore di ogni nodo indica l'ultimo blocco della catena che il nodo considera valida.
% Incluso nel documento con % Fork della blockchain di Bitcoin: evoluzione della rete in cinque passi.
% Il colore di ogni nodo indica l'ultimo blocco della catena che il nodo considera valida.
% Incluso nel documento con \input{immagini/fork.tex}.
\definecolor{forkS}{HTML}{E0E0E0}  % blocco iniziale (stella)
\definecolor{forkX}{HTML}{5B9BD5}  % blocco trovato da X (triangolo)
\definecolor{forkY}{HTML}{ED7D31}  % blocco trovato da Y (triangolo rovesciato)
\definecolor{forkD}{HTML}{2E9E5B}  % nuovo blocco sopra quello di X (rombo)

% posizione dei 12 nodi della rete, uguale in tutti i riquadri
\newcommand{\forkcoordinate}{%
    \coordinate (n1)  at (0.2, 3.0);
    \coordinate (n2)  at (1.4, 3.5);
    \coordinate (n3)  at (2.8, 3.4);
    \coordinate (n4)  at (4.0, 2.9);
    \coordinate (n5)  at (0.7, 2.0);
    \coordinate (n6)  at (1.9, 2.3);
    \coordinate (n7)  at (3.2, 2.0);
    \coordinate (n8)  at (0.1, 0.9);
    \coordinate (n9)  at (1.3, 1.0);
    \coordinate (n10) at (2.6, 0.9);
    \coordinate (n11) at (4.1, 1.3);
    \coordinate (n12) at (1.9, 0.0);
}

% \forkrete{<spostamento>}{<colori dei 12 nodi>}{<didascalia del passo>}
\newcommand{\forkrete}[3]{%
    \begin{scope}[shift={#1}]
        \forkcoordinate
        \foreach \a/\b in {1/2, 2/3, 3/4, 1/5, 2/6, 3/7, 4/7, 4/11, 5/6, 6/7, 5/8, 5/9,
                6/9, 6/10, 7/10, 7/11, 8/9, 9/10, 10/11, 9/12, 10/12}
            \draw[line width=0.8pt] (n\a) -- (n\b);
        \foreach \c [count=\i] in {#2}
            \node[circle, draw, line width=0.6pt, minimum size=0.42cm, fill=\c] at (n\i) {};
        \node[font=\scriptsize\bfseries, above left=4pt] at (n1) {X};
        \node[font=\scriptsize\bfseries, below right=4pt] at (n11) {Y};
        \node[anchor=north, align=center, font=\small] at (2.1, -0.45) {#3};
    \end{scope}
}

\begin{tikzpicture}[
    blocco/.style={draw, fill=white, minimum size=0.85cm, line width=0.8pt},
    simbolo/.style={draw, line width=0.5pt, inner sep=0pt},
    collegamento/.style={-{Latex[length=2mm]}, line width=0.8pt}
]
    \forkrete{(0, 0)}{forkS,forkS,forkS,forkS,forkS,forkS,forkS,forkS,forkS,forkS,forkS,forkS}
        {0. Stato iniziale: tutti i nodi\\hanno la stessa catena}
    \forkrete{(5.8, 0)}{forkX,forkS,forkS,forkS,forkS,forkS,forkS,forkS,forkS,forkS,forkY,forkS}
        {1. Fork: X e Y trovano\\un blocco ciascuno}
    \forkrete{(11.6, 0)}{forkX,forkX,forkY,forkY,forkX,forkX,forkY,forkX,forkX,forkY,forkY,forkY}
        {2. Split della rete}
    \forkrete{(0, -5.4)}{forkD,forkD,forkY,forkY,forkD,forkD,forkD,forkD,forkD,forkD,forkY,forkD}
        {3. Sopra il blocco di X ne viene\\trovato un altro: lo split si espande}
    \forkrete{(5.8, -5.4)}{forkD,forkD,forkD,forkD,forkD,forkD,forkD,forkD,forkD,forkD,forkD,forkD}
        {4. La rete riconverge:\\vince la catena più lunga}

    % legenda: i blocchi e la catena che ne risulta
    \begin{scope}[shift={(11.6, -5.4)}]
        \node[blocco] (bS) at (0.5, 1.7) {};
        \node[simbolo, star, star points=5, star point ratio=2.3, minimum size=0.55cm,
            fill=forkS] at (bS) {};
        \node[blocco] (bX) at (1.9, 2.6) {};
        \node[simbolo, regular polygon, regular polygon sides=3, minimum size=0.6cm,
            fill=forkX] at (bX) {};
        \node[blocco] (bD) at (3.3, 2.6) {};
        \node[simbolo, diamond, minimum size=0.55cm, fill=forkD] at (bD) {};
        \node[blocco, dashed] (bY) at (1.9, 0.8) {};
        \node[simbolo, regular polygon, regular polygon sides=3, shape border rotate=180,
            minimum size=0.6cm, fill=forkY] at (bY) {};
        \draw[collegamento] (bS) -- (bX);
        \draw[collegamento] (bX) -- (bD);
        \draw[collegamento, dashed] (bS) -- (bY);
        \node[font=\small, above=2pt] at ($(bX)!0.5!(bD) + (0, 0.4)$) {catena più lunga};
        \node[font=\small, right=4pt] at (bY.east) {scartato};
        \node[font=\scriptsize, below=1pt] at (bS.south) {iniziale};
        \node[font=\scriptsize, below=1pt] at (bX.south) {di X};
        \node[font=\scriptsize, below=1pt] at (bD.south) {nuovo};
        \node[font=\scriptsize, below=1pt] at (bY.south) {di Y};
        \node[anchor=north, align=center, font=\small] at (2.1, -0.45) {Blocchi e colori\\dei nodi};
    \end{scope}
\end{tikzpicture}
.
\definecolor{forkS}{HTML}{E0E0E0}  % blocco iniziale (stella)
\definecolor{forkX}{HTML}{5B9BD5}  % blocco trovato da X (triangolo)
\definecolor{forkY}{HTML}{ED7D31}  % blocco trovato da Y (triangolo rovesciato)
\definecolor{forkD}{HTML}{2E9E5B}  % nuovo blocco sopra quello di X (rombo)

% posizione dei 12 nodi della rete, uguale in tutti i riquadri
\newcommand{\forkcoordinate}{%
    \coordinate (n1)  at (0.2, 3.0);
    \coordinate (n2)  at (1.4, 3.5);
    \coordinate (n3)  at (2.8, 3.4);
    \coordinate (n4)  at (4.0, 2.9);
    \coordinate (n5)  at (0.7, 2.0);
    \coordinate (n6)  at (1.9, 2.3);
    \coordinate (n7)  at (3.2, 2.0);
    \coordinate (n8)  at (0.1, 0.9);
    \coordinate (n9)  at (1.3, 1.0);
    \coordinate (n10) at (2.6, 0.9);
    \coordinate (n11) at (4.1, 1.3);
    \coordinate (n12) at (1.9, 0.0);
}

% \forkrete{<spostamento>}{<colori dei 12 nodi>}{<didascalia del passo>}
\newcommand{\forkrete}[3]{%
    \begin{scope}[shift={#1}]
        \forkcoordinate
        \foreach \a/\b in {1/2, 2/3, 3/4, 1/5, 2/6, 3/7, 4/7, 4/11, 5/6, 6/7, 5/8, 5/9,
                6/9, 6/10, 7/10, 7/11, 8/9, 9/10, 10/11, 9/12, 10/12}
            \draw[line width=0.8pt] (n\a) -- (n\b);
        \foreach \c [count=\i] in {#2}
            \node[circle, draw, line width=0.6pt, minimum size=0.42cm, fill=\c] at (n\i) {};
        \node[font=\scriptsize\bfseries, above left=4pt] at (n1) {X};
        \node[font=\scriptsize\bfseries, below right=4pt] at (n11) {Y};
        \node[anchor=north, align=center, font=\small] at (2.1, -0.45) {#3};
    \end{scope}
}

\begin{tikzpicture}[
    blocco/.style={draw, fill=white, minimum size=0.85cm, line width=0.8pt},
    simbolo/.style={draw, line width=0.5pt, inner sep=0pt},
    collegamento/.style={-{Latex[length=2mm]}, line width=0.8pt}
]
    \forkrete{(0, 0)}{forkS,forkS,forkS,forkS,forkS,forkS,forkS,forkS,forkS,forkS,forkS,forkS}
        {0. Stato iniziale: tutti i nodi\\hanno la stessa catena}
    \forkrete{(5.8, 0)}{forkX,forkS,forkS,forkS,forkS,forkS,forkS,forkS,forkS,forkS,forkY,forkS}
        {1. Fork: X e Y trovano\\un blocco ciascuno}
    \forkrete{(11.6, 0)}{forkX,forkX,forkY,forkY,forkX,forkX,forkY,forkX,forkX,forkY,forkY,forkY}
        {2. Split della rete}
    \forkrete{(0, -5.4)}{forkD,forkD,forkY,forkY,forkD,forkD,forkD,forkD,forkD,forkD,forkY,forkD}
        {3. Sopra il blocco di X ne viene\\trovato un altro: lo split si espande}
    \forkrete{(5.8, -5.4)}{forkD,forkD,forkD,forkD,forkD,forkD,forkD,forkD,forkD,forkD,forkD,forkD}
        {4. La rete riconverge:\\vince la catena più lunga}

    % legenda: i blocchi e la catena che ne risulta
    \begin{scope}[shift={(11.6, -5.4)}]
        \node[blocco] (bS) at (0.5, 1.7) {};
        \node[simbolo, star, star points=5, star point ratio=2.3, minimum size=0.55cm,
            fill=forkS] at (bS) {};
        \node[blocco] (bX) at (1.9, 2.6) {};
        \node[simbolo, regular polygon, regular polygon sides=3, minimum size=0.6cm,
            fill=forkX] at (bX) {};
        \node[blocco] (bD) at (3.3, 2.6) {};
        \node[simbolo, diamond, minimum size=0.55cm, fill=forkD] at (bD) {};
        \node[blocco, dashed] (bY) at (1.9, 0.8) {};
        \node[simbolo, regular polygon, regular polygon sides=3, shape border rotate=180,
            minimum size=0.6cm, fill=forkY] at (bY) {};
        \draw[collegamento] (bS) -- (bX);
        \draw[collegamento] (bX) -- (bD);
        \draw[collegamento, dashed] (bS) -- (bY);
        \node[font=\small, above=2pt] at ($(bX)!0.5!(bD) + (0, 0.4)$) {catena più lunga};
        \node[font=\small, right=4pt] at (bY.east) {scartato};
        \node[font=\scriptsize, below=1pt] at (bS.south) {iniziale};
        \node[font=\scriptsize, below=1pt] at (bX.south) {di X};
        \node[font=\scriptsize, below=1pt] at (bD.south) {nuovo};
        \node[font=\scriptsize, below=1pt] at (bY.south) {di Y};
        \node[anchor=north, align=center, font=\small] at (2.1, -0.45) {Blocchi e colori\\dei nodi};
    \end{scope}
\end{tikzpicture}
.
\definecolor{forkS}{HTML}{E0E0E0}  % blocco iniziale (stella)
\definecolor{forkX}{HTML}{5B9BD5}  % blocco trovato da X (triangolo)
\definecolor{forkY}{HTML}{ED7D31}  % blocco trovato da Y (triangolo rovesciato)
\definecolor{forkD}{HTML}{2E9E5B}  % nuovo blocco sopra quello di X (rombo)

% posizione dei 12 nodi della rete, uguale in tutti i riquadri
\newcommand{\forkcoordinate}{%
    \coordinate (n1)  at (0.2, 3.0);
    \coordinate (n2)  at (1.4, 3.5);
    \coordinate (n3)  at (2.8, 3.4);
    \coordinate (n4)  at (4.0, 2.9);
    \coordinate (n5)  at (0.7, 2.0);
    \coordinate (n6)  at (1.9, 2.3);
    \coordinate (n7)  at (3.2, 2.0);
    \coordinate (n8)  at (0.1, 0.9);
    \coordinate (n9)  at (1.3, 1.0);
    \coordinate (n10) at (2.6, 0.9);
    \coordinate (n11) at (4.1, 1.3);
    \coordinate (n12) at (1.9, 0.0);
}

% \forkrete{<spostamento>}{<colori dei 12 nodi>}{<didascalia del passo>}
\newcommand{\forkrete}[3]{%
    \begin{scope}[shift={#1}]
        \forkcoordinate
        \foreach \a/\b in {1/2, 2/3, 3/4, 1/5, 2/6, 3/7, 4/7, 4/11, 5/6, 6/7, 5/8, 5/9,
                6/9, 6/10, 7/10, 7/11, 8/9, 9/10, 10/11, 9/12, 10/12}
            \draw[line width=0.8pt] (n\a) -- (n\b);
        \foreach \c [count=\i] in {#2}
            \node[circle, draw, line width=0.6pt, minimum size=0.42cm, fill=\c] at (n\i) {};
        \node[font=\scriptsize\bfseries, above left=4pt] at (n1) {X};
        \node[font=\scriptsize\bfseries, below right=4pt] at (n11) {Y};
        \node[anchor=north, align=center, font=\small] at (2.1, -0.45) {#3};
    \end{scope}
}

\begin{tikzpicture}[
    blocco/.style={draw, fill=white, minimum size=0.85cm, line width=0.8pt},
    simbolo/.style={draw, line width=0.5pt, inner sep=0pt},
    collegamento/.style={-{Latex[length=2mm]}, line width=0.8pt}
]
    \forkrete{(0, 0)}{forkS,forkS,forkS,forkS,forkS,forkS,forkS,forkS,forkS,forkS,forkS,forkS}
        {0. Stato iniziale: tutti i nodi\\hanno la stessa catena}
    \forkrete{(5.8, 0)}{forkX,forkS,forkS,forkS,forkS,forkS,forkS,forkS,forkS,forkS,forkY,forkS}
        {1. Fork: X e Y trovano\\un blocco ciascuno}
    \forkrete{(11.6, 0)}{forkX,forkX,forkY,forkY,forkX,forkX,forkY,forkX,forkX,forkY,forkY,forkY}
        {2. Split della rete}
    \forkrete{(0, -5.4)}{forkD,forkD,forkY,forkY,forkD,forkD,forkD,forkD,forkD,forkD,forkY,forkD}
        {3. Sopra il blocco di X ne viene\\trovato un altro: lo split si espande}
    \forkrete{(5.8, -5.4)}{forkD,forkD,forkD,forkD,forkD,forkD,forkD,forkD,forkD,forkD,forkD,forkD}
        {4. La rete riconverge:\\vince la catena più lunga}

    % legenda: i blocchi e la catena che ne risulta
    \begin{scope}[shift={(11.6, -5.4)}]
        \node[blocco] (bS) at (0.5, 1.7) {};
        \node[simbolo, star, star points=5, star point ratio=2.3, minimum size=0.55cm,
            fill=forkS] at (bS) {};
        \node[blocco] (bX) at (1.9, 2.6) {};
        \node[simbolo, regular polygon, regular polygon sides=3, minimum size=0.6cm,
            fill=forkX] at (bX) {};
        \node[blocco] (bD) at (3.3, 2.6) {};
        \node[simbolo, diamond, minimum size=0.55cm, fill=forkD] at (bD) {};
        \node[blocco, dashed] (bY) at (1.9, 0.8) {};
        \node[simbolo, regular polygon, regular polygon sides=3, shape border rotate=180,
            minimum size=0.6cm, fill=forkY] at (bY) {};
        \draw[collegamento] (bS) -- (bX);
        \draw[collegamento] (bX) -- (bD);
        \draw[collegamento, dashed] (bS) -- (bY);
        \node[font=\small, above=2pt] at ($(bX)!0.5!(bD) + (0, 0.4)$) {catena più lunga};
        \node[font=\small, right=4pt] at (bY.east) {scartato};
        \node[font=\scriptsize, below=1pt] at (bS.south) {iniziale};
        \node[font=\scriptsize, below=1pt] at (bX.south) {di X};
        \node[font=\scriptsize, below=1pt] at (bD.south) {nuovo};
        \node[font=\scriptsize, below=1pt] at (bY.south) {di Y};
        \node[anchor=north, align=center, font=\small] at (2.1, -0.45) {Blocchi e colori\\dei nodi};
    \end{scope}
\end{tikzpicture}
.
\definecolor{forkS}{HTML}{E0E0E0}  % blocco iniziale (stella)
\definecolor{forkX}{HTML}{5B9BD5}  % blocco trovato da X (triangolo)
\definecolor{forkY}{HTML}{ED7D31}  % blocco trovato da Y (triangolo rovesciato)
\definecolor{forkD}{HTML}{2E9E5B}  % nuovo blocco sopra quello di X (rombo)

% posizione dei 12 nodi della rete, uguale in tutti i riquadri
\newcommand{\forkcoordinate}{%
    \coordinate (n1)  at (0.2, 3.0);
    \coordinate (n2)  at (1.4, 3.5);
    \coordinate (n3)  at (2.8, 3.4);
    \coordinate (n4)  at (4.0, 2.9);
    \coordinate (n5)  at (0.7, 2.0);
    \coordinate (n6)  at (1.9, 2.3);
    \coordinate (n7)  at (3.2, 2.0);
    \coordinate (n8)  at (0.1, 0.9);
    \coordinate (n9)  at (1.3, 1.0);
    \coordinate (n10) at (2.6, 0.9);
    \coordinate (n11) at (4.1, 1.3);
    \coordinate (n12) at (1.9, 0.0);
}

% \forkrete{<spostamento>}{<colori dei 12 nodi>}{<didascalia del passo>}
\newcommand{\forkrete}[3]{%
    \begin{scope}[shift={#1}]
        \forkcoordinate
        \foreach \a/\b in {1/2, 2/3, 3/4, 1/5, 2/6, 3/7, 4/7, 4/11, 5/6, 6/7, 5/8, 5/9,
                6/9, 6/10, 7/10, 7/11, 8/9, 9/10, 10/11, 9/12, 10/12}
            \draw[line width=0.8pt] (n\a) -- (n\b);
        \foreach \c [count=\i] in {#2}
            \node[circle, draw, line width=0.6pt, minimum size=0.42cm, fill=\c] at (n\i) {};
        \node[font=\scriptsize\bfseries, above left=4pt] at (n1) {X};
        \node[font=\scriptsize\bfseries, below right=4pt] at (n11) {Y};
        \node[anchor=north, align=center, font=\small] at (2.1, -0.45) {#3};
    \end{scope}
}

\begin{tikzpicture}[
    blocco/.style={draw, fill=white, minimum size=0.85cm, line width=0.8pt},
    simbolo/.style={draw, line width=0.5pt, inner sep=0pt},
    collegamento/.style={-{Latex[length=2mm]}, line width=0.8pt}
]
    \forkrete{(0, 0)}{forkS,forkS,forkS,forkS,forkS,forkS,forkS,forkS,forkS,forkS,forkS,forkS}
        {0. Stato iniziale: tutti i nodi\\hanno la stessa catena}
    \forkrete{(5.8, 0)}{forkX,forkS,forkS,forkS,forkS,forkS,forkS,forkS,forkS,forkS,forkY,forkS}
        {1. Fork: X e Y trovano\\un blocco ciascuno}
    \forkrete{(11.6, 0)}{forkX,forkX,forkY,forkY,forkX,forkX,forkY,forkX,forkX,forkY,forkY,forkY}
        {2. Split della rete}
    \forkrete{(0, -5.4)}{forkD,forkD,forkY,forkY,forkD,forkD,forkD,forkD,forkD,forkD,forkY,forkD}
        {3. Sopra il blocco di X ne viene\\trovato un altro: lo split si espande}
    \forkrete{(5.8, -5.4)}{forkD,forkD,forkD,forkD,forkD,forkD,forkD,forkD,forkD,forkD,forkD,forkD}
        {4. La rete riconverge:\\vince la catena più lunga}

    % legenda: i blocchi e la catena che ne risulta
    \begin{scope}[shift={(11.6, -5.4)}]
        \node[blocco] (bS) at (0.5, 1.7) {};
        \node[simbolo, star, star points=5, star point ratio=2.3, minimum size=0.55cm,
            fill=forkS] at (bS) {};
        \node[blocco] (bX) at (1.9, 2.6) {};
        \node[simbolo, regular polygon, regular polygon sides=3, minimum size=0.6cm,
            fill=forkX] at (bX) {};
        \node[blocco] (bD) at (3.3, 2.6) {};
        \node[simbolo, diamond, minimum size=0.55cm, fill=forkD] at (bD) {};
        \node[blocco, dashed] (bY) at (1.9, 0.8) {};
        \node[simbolo, regular polygon, regular polygon sides=3, shape border rotate=180,
            minimum size=0.6cm, fill=forkY] at (bY) {};
        \draw[collegamento] (bS) -- (bX);
        \draw[collegamento] (bX) -- (bD);
        \draw[collegamento, dashed] (bS) -- (bY);
        \node[font=\small, above=2pt] at ($(bX)!0.5!(bD) + (0, 0.4)$) {catena più lunga};
        \node[font=\small, right=4pt] at (bY.east) {scartato};
        \node[font=\scriptsize, below=1pt] at (bS.south) {iniziale};
        \node[font=\scriptsize, below=1pt] at (bX.south) {di X};
        \node[font=\scriptsize, below=1pt] at (bD.south) {nuovo};
        \node[font=\scriptsize, below=1pt] at (bY.south) {di Y};
        \node[anchor=north, align=center, font=\small] at (2.1, -0.45) {Blocchi e colori\\dei nodi};
    \end{scope}
\end{tikzpicture}
